Portfolio risk assessment ordinarily relies on reliable estimates of
cross-asset return covariances, which are difficult to obtain in short,
high-dimensional panels.
We show that firm-level distribution-valued characteristics can instead provide
one-sided certificates of portfolio risk.
Under maintained links from characteristics to systematic exposures and from
exposures to returns, multi-firm Wasserstein-2 dispersion yields a sharp upper
bound on systematic portfolio variance and a corresponding bound for
standardized returns.
A weighted pairwise relaxation produces an objective that is convex under a
checkable condition and requires marginal volatility scales but no cross-asset
return covariances.
With zero firm-specific slack, the common-map scale changes the certified
variance reduction but not the normalized allocation, which depends only on
observed information geometry.
In a \ppnum[0]{n-tickers}-firm panel from 2018--2022, only
\ppnum[2]{news-only-percentile-uniform-cap-15-pct}\%--
\ppnum[2]{news-only-percentile-concentrated-cap-15-pct}\% of portfolios across
four prespecified capped reference populations have in-sample variance no
greater than an allocation constructed from Qwen3-Embedding-8B news
representations; the corresponding share for equal risk weighting is
\ppnum[1]{equal-percentile-concentrated-cap-15-pct}\%--
\ppnum[1]{equal-percentile-uniform-cap-12p5-pct}\%.
The lower in-sample variance ranking relative to equal risk also appears across
the reported frozen language-model representations.
The framework therefore converts distribution-valued firm information into a
coherent risk bound and an implementable allocation rule constructed without
cross-asset return covariances.
